\documentclass[11pt]{article}
\usepackage[margin=1in]{geometry}
\usepackage{booktabs}
\title{VTA Co-expression Network --- Interpretation}
\author{GSE233866, mouse midbrain DA neurons, healthy baseline}
\date{}

\begin{document}
\maketitle

VTA's soft power scan independently selects $\beta=5$ ($R^2=0.850$), matching
SNc's choice by coincidence of the data rather than by assumption --- each
network's power was fit separately. Clustering on only 832 pooled VTA cells
(versus SNc's 2{,}226) yields 39 modules, more than four times SNc's count.
This is best read as a sample-size artifact rather than a genuine difference
in how organized VTA's co-expression structure is: fewer cells produce
noisier gene-gene correlations, which tends to fragment clustering into many
small modules rather than a few large, stable ones. Many of VTA's 39 modules
are under 100 genes.

CACNA1D is assigned to the \textbf{brown} module with $\mathrm{kME}=0.296$ ---
weak-to-moderate, and \emph{lower} than its SNc value (0.442), the opposite
of what this document previously (incorrectly) reported.
\emph{(Correction, 2026-08-13: this document previously reported
$\mathrm{kME}=0.094$ and described it as roughly double the SNc value;
that was misread from the wrong \texttt{kME\_*} column in
\texttt{modules.csv}. 0.296 is \texttt{kME\_brown} for Cacna1d, re-verified
directly against the source file --- SNc, not VTA, has the stronger
connection.)}

Brown is significantly enriched for \emph{calcium ion transmembrane import
into cytosol} ($p_{\mathrm{adj}} = 7.8\times10^{-3}$, GO term, unaffected by
the KEGG-library fix applied elsewhere in this project) --- directly
on-topic for a calcium-channel gene, though a weaker enrichment signal than
SNc turquoise's dopaminergic-synapse hit. Both regions place CACNA1D in a
functionally sensible, moderately-central module --- SNc's connection is the
stronger of the two.

\end{document}
